% Draft Statutory Text. Built by the future pass from the instructions below. Do
% not process now. This is the operative bill: real, enactable statutory language
% in proper legislative form. The model for wholly new Physical AI statutory text
% is the template's new Subpart J; the validation thresholds come from the
% VVUQ-02 gate spec. Use the section symbol § throughout and the lettered,
% numbered, roman subparagraph hierarchy (a)(1)(i).

[Draft the statutory text as enactable legislative language, organised into
numbered sections with the (a)(1)(i) subparagraph hierarchy, in the style and tone
of the template's new Subpart J. Build each section from the sources named below;
keep any threshold tables in the left-aligned ragged-right format and reference
them from the text rather than restating the numbers in prose. Do not narrate the
sources; convert them into statute. Suggested sections follow.]

[Section 1, Short title. "VVUQ Physical AI Oncology Trial Bill." One clause.]

[Section 2, Definitions. Incorporate by reference the defined terms set out in the
bill's definitions section, listing the operative terms (Physical AI system,
robot-patient interaction, VVUQ gate, verification fraction, USL, PSL, high-risk
artificial intelligence, AI-enabled device software function) without re-defining
them. Cross-reference \texttt{sections/definitions.tex}.]

[Section 3, The verification-before-generation requirement (the core of the bill).
Draft the operative mandate: no robot-patient interaction code may be generated or
executed in a Physical AI oncology clinical trial unless an automated VVUQ
verification process, bound to named external standards, has first cleared for
that interaction, and the cleared verification record has been documented and
attested. Make the priority explicit, that the verification process is held to a
higher standard than, and precedes, the generated code. Draw the structure from
the template's \texttt{chunk\_10\_subpart\_j\_physical\_ai\_references.md}
(§ 312.402 Validation Requirements and § 312.400 Purpose and Scope) and the gate
mechanics from \texttt{VVUQ-02/codegen/docs/vvuq\_gate\_spec.md} and
\texttt{vvuq\_methodology.md}. Cite \texttt{kawchak2026vvuq01},
\texttt{kawchak2026vvuq02}, \texttt{asmevv40}, and \texttt{fda2023credibility}.]

[Section 4, Validation requirements and gate thresholds. Codify, in subparagraph
form and one threshold table, the gate requirements a robot-patient interaction
must clear: a verification fraction equal to 1.0; a validation agreement against an
independent reference; a coefficient-of-variation dispersion bound across repeated
seeded runs; and a hard predicate for immediate-catastrophe interactions (vascular
no-fly, shared-operating-room collision, fault and e-stop). Source the ten-gate
threshold structure from \texttt{vvuq\_gate\_spec.md} and bind each requirement to
its external standard via \texttt{standards\_map.md}. Cite \texttt{iso14971},
\texttt{iec8060127}, \texttt{isots15066}, \texttt{nasastd7009a}, \texttt{ul4600},
and \texttt{ieee7009}.]

[Section 5, Readiness gates. Require that the site clear the PSL pass/fail gate and
the robot clear the minimum USL score for the procedure type (surgical at least
7.0; therapeutic positioning and diagnostic needle at least 6.0; rehabilitative at
least 4.0; companion monitoring at least 3.0) before enrollment, and require Phase
0 simulation validation. Source from the national-platform
\texttt{psl\_usl\_standards.tex} and the template's § 312.401 classification and
§ 312.402 validation. Cite \texttt{kawchak2026national}.]

[Section 6, Documentation and attestation. Require the algorithm documentation and
the compliance statement and attestation set out in the bill's documentation and
attestation sections, by cross-reference, and require the verification record to
be filed before generation. Cross-reference
\texttt{sections/algorithm\_documentation.tex} and
\texttt{sections/attestations\_compliance.tex}, and cite \texttt{tx-sb1822},
\texttt{ny-a9149}, and \texttt{ca-sb1120}.]

[Section 7, Cybersecurity, human oversight, and lifecycle. Codify the cybersecurity
requirements, the recorded human-oversight and hand-back-to-human escalation, and
the lifecycle and decommissioning duties, drawn from the template's § 312.403
(cybersecurity), § 312.404 (human oversight), and § 312.405 (lifecycle
management). Tie the electronic-records integrity to \texttt{cfr-part11} and the
hash-chained audit trail to the national MCP-PAI standard. Cite
\texttt{kawchak2026mcp} and \texttt{iec60601}.]

[Section 8, Nondiscrimination. Require that algorithms and their training data
minimize the risk of bias on protected characteristics and adhere to evidence-
based clinical guidelines, consistent with \texttt{usc-titlevi},
\texttt{usc-ada}, and \texttt{ny-a9149}.]

[Section 9, Enforcement and effective date. Cross-reference the enforcement and
fiscal section and the effective-date clause in the conclusions section; do not
restate them. Cross-reference \texttt{sections/implementation\_enforcement.tex} and
\texttt{sections/conclusions.tex}.]

[Add a closing authority-and-source line in the template's bracketed-citation
style (modelled on \texttt{chunk\_01\_front\_matter.md}), making clear the bill
adapts public-domain federal material and is not endorsed by CFR, ICH, or FDA.]
